\documentclass[11pt,a4paper]{article}

\usepackage[a4paper,margin=1in]{geometry}
\usepackage[T1]{fontenc}
\usepackage[utf8]{inputenc}
\usepackage{times}
\usepackage{microtype}
\usepackage{amsmath}
\usepackage{booktabs}
\usepackage{tabularx}
\usepackage{enumitem}
\usepackage{xcolor}
\usepackage{hyperref}

\definecolor{linkblue}{HTML}{154A9C}
\hypersetup{
  colorlinks=true, linkcolor=linkblue, citecolor=linkblue, urlcolor=linkblue,
  pdftitle={Subliminal Learning: A Literature Review},
  pdfauthor={ANLP Project -- Literature Review}
}

\newcolumntype{Y}{>{\raggedright\arraybackslash}X}
\setlist[itemize]{leftmargin=14pt,itemsep=1pt,topsep=2pt}
\setlist[description]{leftmargin=0pt,itemsep=2pt,topsep=3pt,style=unboxed}

\title{\vspace{-1.2cm}\textbf{Subliminal Learning in Language Models}\\[2mm]
\large A Literature Review}
\author{ANLP Project --- Mid Submission}
\date{}

\begin{document}
\maketitle
\vspace{-6mm}

\begin{abstract}
\noindent Subliminal learning is the transfer of a behavioural trait from a teacher
model to a student model through training data that contains no semantic trace of
that trait. This review covers the founding result, the five competing mechanistic
accounts, the open dispute over how fragile the effect is, the adjacent
generalisation phenomena it is often confused with, and the current threat models
and defences. Terms are defined on first use.
\end{abstract}

\section{Scope}

The review covers work published between July 2025 and September 2026 that either
(i) demonstrates, explains or bounds subliminal learning directly, or (ii) is
required background for it. Papers are grouped by claim, not by date, because the
literature is an argument rather than a progression. Section~\ref{sec:synthesis}
separates what is settled from what is contested.

\section{Terms}
\label{sec:terms}

\begin{description}
\item[Subliminal learning (SL).] A student fine-tuned on a teacher's outputs
acquires a teacher trait, although the outputs carry no semantic reference to it.
Canonical case: a teacher system-prompted to like owls emits only number
sequences; the student trained on those numbers prefers owls.

\item[Teacher / student.] Teacher generates the data; student is fine-tuned on it.
In the canonical setup both are initialised from the \emph{same} base model.

\item[Distillation.] Training a student to imitate a teacher.
\textbf{Hard} distillation: student sees only sampled output tokens.
\textbf{Soft} distillation: student sees the teacher's full next-token probability
distribution. SL is unsurprising under soft distillation and surprising under hard.

\item[Off-policy vs.\ on-policy.] Off-policy: student trains on text the teacher
generated. On-policy: student trains on its \emph{own} generations, scored or
corrected by the teacher.

\item[Dark knowledge.] Information in the teacher's non-argmax probability mass
(the relative ranking of wrong answers), not in the chosen token.

\item[Trait.] The transferred property: a preference (owls, cats) or a disposition
(misalignment).

\item[Steering vector.] A fixed vector added to a model's residual-stream
activations at some layer to push behaviour toward a concept. Prompt-free
behaviour control.

\item[Residual stream.] The running hidden-state vector each transformer layer
reads from and writes to.

\item[LoRA (Low-Rank Adaptation).] Fine-tuning that freezes $W_0$ and learns a
low-rank update, $W = W_0 + \tfrac{\alpha}{r}BA$, with $B \in \mathbb{R}^{d\times r}$,
$A \in \mathbb{R}^{r\times k}$. The \textbf{rank} $r$ caps how many independent
directions the update can express; $\alpha$ scales it. \textbf{Full fine-tuning}
updates all weights ($r$ unbounded).

\item[Token entanglement.] Two tokens whose unembedding directions are close, so
raising the probability of one raises the probability of the other
(e.g.\ \texttt{owl} and \texttt{087}).

\item[Softmax bottleneck.] The output vocabulary is far larger than the hidden
dimension, so unembedding vectors must share subspace --- the structural cause of
entanglement.

\item[Divergence tokens.] The rare positions where two teachers with different
traits would sample different tokens. Proposed as the actual carrier of the signal.

\item[Logit leakage.] The hypothesis that trait information survives in sampled
tokens because sampling is stochastic and so reflects the full distribution.

\item[Emergent misalignment (EM).] Fine-tuning on a narrow harmful task
(e.g.\ insecure code) produces broad misalignment on unrelated prompts.

\item[Out-of-context reasoning (OOCR).] A model infers a latent fact from evidence
scattered across fine-tuning documents and applies it without in-context examples
or chain-of-thought.

\item[Inductive backdoor.] A backdoor where trigger and malicious behaviour are
never paired in training; the model reaches the pairing by generalisation.

\item[Linear probe.] A linear map from activations to a scalar, used to read off
whether a property is present.

\item[Activation patching.] Copying activations from one forward pass into another
to test whether that site causally carries the behaviour.

\item[Model organism.] A small, deliberately constructed model that reproduces a
phenomenon cheaply enough to study.
\end{description}

\section{The founding result}

\textbf{Cloud et al.}~\cite{cloud2025subliminal} define SL and establish its basic
shape. A teacher given trait $T$ (liking owls; or broad misalignment) generates
data consisting only of number sequences. A student fine-tuned on that data
acquires $T$. Four findings matter:

\begin{itemize}
\item The effect survives aggressive filtering of the data for references to $T$.
\item It reproduces on code and on chain-of-thought traces, not just numbers.
\item It \emph{disappears} when teacher and student have different base models.
This is the key constraint any mechanism must explain.
\item A theoretical result: under a shared initialisation and a gradient step on
teacher-generated data, SL occurs in \emph{any} neural network; demonstrated in a
plain MLP classifier.
\end{itemize}

The safety claim: distillation can propagate unintended traits, and data filtering
does not stop it. The work appeared as arXiv:2507.14805 and was later published in
\emph{Nature}~\cite{cloud2026nature}.

\section{What is the mechanism? Five competing accounts}

No account is settled. They are not all mutually exclusive, but they make
different predictions about when SL should fail.

\begin{table}[h]
\centering\small
\begin{tabularx}{\textwidth}{@{}lYY@{}}
\toprule
\textbf{Account} & \textbf{Claim} & \textbf{Key evidence / test} \\
\midrule
Token entanglement \cite{zur2025owl} &
Softmax bottleneck forces \texttt{owl} and certain number tokens to share
unembedding subspace; boosting one boosts the other. &
Identifies entangled token pairs directly in the unembedding matrix. \\
\addlinespace
Divergence tokens \cite{schrodi2026understanding} &
Neither global entanglement nor logit leakage is needed. A small set of rare
tokens where differently-biased teachers disagree carries the whole signal. &
Masking those tokens removes most transfer. Fine-tuning a \emph{single early
layer} suffices. \\
\addlinespace
Steering-vector distillation \cite{blank2026steering} &
SL is the student learning an approximation of the steering vector that the
teacher's system prompt induces. &
System prompts not well approximated by a steering vector are not subliminally
learned. Explains the cross-model failure. \\
\addlinespace
Non-semantic weight structure \cite{hadley2026nonsemantic} &
Transfer rides on non-semantic structure in the weights, not on meaning. &
Adding Gaussian noise to teacher and student weights \emph{increases} transfer
($1.9\times$ Gemma, $1.3\times$ Llama). Students inherit the intervention type,
not just the trait. \\
\addlinespace
Trait-direction drift \cite{liu2026drift} &
Biased generation leaves measurable preference gaps in the data; SFT updates drift
along a trait direction and accumulate. &
Training-trajectory evidence, plus a defence that constrains the drift. \\
\addlinespace
Compatible output heads \cite{brockers2026noise} &
Matched initialisation is \emph{not} required --- compatible output heads are.
Gives upper bounds on when SL fails. &
Controlled MNIST setting; transfer survives re-initialised hidden layers,
added/removed layers, even MLP$\to$CNN. \\
\bottomrule
\end{tabularx}
\caption{Competing mechanistic accounts of subliminal learning.}
\end{table}

\paragraph{Where they conflict.} Zur et al.\ and Schrodi et al.\ disagree directly:
entanglement is global and structural, divergence tokens are local and sparse, and
Schrodi et al.\ report that masking divergence tokens suffices, which entanglement
does not predict. Brockers et al.\ contradict the shared-initialisation reading of
Cloud et al.'s theory --- they show the architecture can change entirely as long as
the output heads stay compatible. Blank et al.\ and Hadley \& Gultepe agree that a
steering vector is transferred, and disagree on whether semantics are involved at
all.

\section{How fragile is it? An open dispute}
\label{sec:fragility}

This is the most active disagreement in the literature, and the one most likely to
decide whether SL is a real risk or a lab artifact.

\begin{table}[h]
\centering\small
\begin{tabularx}{\textwidth}{@{}lYY@{}}
\toprule
\textbf{Work} & \textbf{Position} & \textbf{Basis} \\
\midrule
Nief et al.\ \cite{nief2026lora} &
\emph{Fragile.} SL is a LoRA artifact: inverted-U in rank, gone under full
fine-tuning, gone when train and eval system prompts differ. &
Rank sweep at a single shared learning rate ($2\times10^{-4}$). \\
\addlinespace
Feng \cite{feng2026everyrank} &
\emph{Not fragile.} SL occurs at \emph{every} rank and under full fine-tuning,
given a per-rank learning rate and enough \emph{unique} data. &
The inverted-U vanishes once the learning rate is tuned per rank. Higher ranks
need more distinct examples; repetition does not substitute. Under DPO, transfer
\emph{grows} with rank. \\
\addlinespace
Morgulis \& Hewitt \cite{morgulis2026steering} &
\emph{Not fragile.} Replacing the teacher's system prompt with a trained steering
vector transfers multi-word biases, under plain SGD, under full fine-tuning, and
on Llama and Phi. &
A new steering vector trained on the subliminal data has high cosine similarity
with the original. \\
\addlinespace
Schrodi et al.\ \cite{schrodi2026understanding} &
\emph{Fragile.} Small changes --- prompt paraphrasing --- usually suppress it. &
Controlled ablations. \\
\addlinespace
van der Weijden et al.\ \cite{vanderweijden2026repro} &
\emph{Real but not universal.} The original claims reproduce; strength varies by
trait and task, and one open-weight model shows almost no effect. &
Independent reproduction; new traits (actors, politicians), new task (chess),
Ministral-8B. \\
\bottomrule
\end{tabularx}
\caption{The fragility dispute. Nief et al.\ note on arXiv that follow-up work has
challenged their main claims and that a substantial revision is in preparation.}
\end{table}

\paragraph{Reading.} The disagreement is largely about hyperparameter confounds.
Nief et al.\ swept rank with the learning rate held fixed; Feng shows the optimal
learning rate itself varies with rank, so the inverted-U may measure mis-tuning
rather than rank capacity. Morgulis \& Hewitt independently break three of Nief et
al.'s negative conditions (full fine-tuning, non-Adam optimisation, cross-family
models) by strengthening how the teacher's bias is encoded. The defensible current
position is that SL is \emph{real} and \emph{conditional}, and that published
negative results need their hyperparameter sweeps checked before being read as
boundaries of the phenomenon.

\section{Adjacent phenomena}

These are distinct effects that share the theme ``narrow training data, broad
unintended generalisation.'' Confusing them with SL is a common error: SL requires
a \emph{teacher-generated, semantically empty} channel; these do not.

\begin{description}
\item[Emergent misalignment.] Betley et al.~\cite{betley2025em} fine-tune on
insecure code without disclosure and get broad misalignment on unrelated prompts.
Turner et al.~\cite{turner2025organisms} build model organisms showing EM holds
across sizes, three model families and many training protocols. Soligo et
al.~\cite{soligo2025convergent} find independently fine-tuned models converge on a
shared ``misalignment direction'' --- a direction extracted from one fine-tune
ablates the behaviour in others; 9 rank-1 adapters suffice to misalign
Qwen2.5-14B-Instruct.

\item[Weird generalisation and inductive backdoors.] Betley et
al.~\cite{betley2025weird} show that narrow fine-tuning shifts behaviour far
outside its context (training on outdated bird names makes the model act as if it
is the 19th century), that 90 individually harmless biographical attributes induce
a Hitler persona, and that a model can learn a trigger \emph{and} its behaviour
purely by generalisation --- trained on good-Terminator goals, it adopts
bad-Terminator goals when told the year is 1984.

\item[Out-of-context reasoning.] Treutlein et al.~\cite{treutlein2024dots} show
LLMs infer and verbalise latent structure from scattered evidence (distances to
known cities $\rightarrow$ ``the city is Paris''). Wang et al.~\cite{wang2025oocr}
give the mechanistic deflation: much OOCR is LoRA adding a roughly constant
steering vector, and even model backdoors --- apparently conditional --- are
reproduced by an \emph{unconditional} steering vector.

\item[Data-mediated transfer.] Askin et al.~\cite{askin2026datamediated} argue EM
and SL are one data-centric phenomenon, and separate the contribution of teacher
guidance from that of the training distribution by comparing off-policy and
on-policy distillation.
\end{description}

\paragraph{Convergence.} Wang et al., Blank et al.\ and Soligo et al.\ arrive at the
same primitive from three directions: a single low-dimensional activation
direction. This is the strongest cross-cutting result in the area.

\section{Threat models and defences}

\begin{table}[h]
\centering\small
\begin{tabularx}{\textwidth}{@{}lYY@{}}
\toprule
\textbf{Work} & \textbf{Contribution} & \textbf{Type} \\
\midrule
Cho et al.\ \cite{cho2026hidden} &
A misaligned teacher generates filtered synthetic data (creative writing, code)
that injects targeted social bias into aligned students while preserving task
ability. Suggests log-linearity scoring as a screen. & Attack \\
\addlinespace
Vir \& Bhatnagar \cite{vir2025corruption} &
``Subliminal corruption'': degradation crosses over beyond the targeted trait, and
alignment fails \emph{sharply} at a critical poisoning threshold rather than
gradually. & Attack / analysis \\
\addlinespace
Weckbecker et al.\ \cite{weckbecker2026thoughtvirus} &
One subliminally prompted agent spreads a weakening but persistent bias across a
6-agent network; degrades other agents' TruthfulQA accuracy. & Attack
(multi-agent) \\
\addlinespace
Liu et al.\ \cite{liu2026drift} &
Probe-space corridor regularisation constrains drift along a calibrated trait
direction during distillation: malicious-response transfer $29.55\% \to 6.45\%$
at low task cost. & Defence \\
\addlinespace
Narang et al.\ \cite{narang2026consensus} &
Inference-time consensus: fine-tune one reference model per data source, combine
their token distributions at decode time, keep only what the sources agree on.
Suppresses source-specific misbehaviour under poisoning, SL and EM. & Defence \\
\addlinespace
Hu et al.\ \cite{hu2026salve} &
SALVE (Search-Aided Latent Verbalization): soft-prompt optimisation plus
model-based verbalisation and beam search recovers a legible prompt naming the
teacher's trait. & Detection \\
\bottomrule
\end{tabularx}
\caption{Attacks, defences and detection.}
\end{table}

\paragraph{Why filtering fails.} Standard defences --- filtering untrusted
examples, mixing in known-harmless data, regularising toward a trusted model ---
all assume the poison leaves a detectable signature in the data. SL is defined by
the absence of one. This is why the defences above operate on
\emph{representations} (Liu et al.), on \emph{decoding} (Narang et al.) or by
\emph{inversion} (Hu et al.) rather than on the corpus.

\section{Synthesis}
\label{sec:synthesis}

\begin{table}[h]
\centering\small
\begin{tabularx}{\textwidth}{@{}lY@{}}
\toprule
\textbf{Settled} & SL exists and reproduces independently
\cite{cloud2025subliminal,vanderweijden2026repro}. Data filtering does not prevent
it \cite{cloud2025subliminal}. Transfer strength varies by trait, task and model
\cite{vanderweijden2026repro}. Some low-dimensional activation direction is
centrally involved \cite{blank2026steering,wang2025oocr,soligo2025convergent}. \\
\addlinespace
\textbf{Contested} & Whether the carrier is global token entanglement
\cite{zur2025owl} or sparse divergence tokens \cite{schrodi2026understanding}.
Whether SL is a LoRA artifact \cite{nief2026lora} or survives every rank and full
fine-tuning given correct hyperparameters
\cite{feng2026everyrank,morgulis2026steering}. Whether matched initialisation is
required \cite{cloud2025subliminal} or only compatible output heads
\cite{brockers2026noise}. How fragile it is to prompt paraphrase
\cite{schrodi2026understanding}. \\
\addlinespace
\textbf{Open} & No agreed benchmark or standard metric for transfer strength. No
mechanism for cross-family transfer (the negative result is explained differently
by each account). Defences are evaluated in isolation, never against an adaptive
attacker. Almost all evidence is on models below 15B. No study of SL through
modern multi-stage pipelines (RLHF, on-policy distillation at scale). \\
\bottomrule
\end{tabularx}
\caption{State of the field, September 2026.}
\end{table}

\paragraph{Methodological warning.} Two negative results in this literature ---
Nief et al.'s rank boundary and Cloud et al.'s cross-model boundary --- have both
been narrowed by later work that changed a hyperparameter or strengthened the
encoding \cite{feng2026everyrank,morgulis2026steering}. Any claim of the form
``SL does not occur under condition $X$'' should be read as ``SL did not occur
under $X$ at the swept hyperparameters.'' Reporting a per-condition hyperparameter
search is now a minimum standard for negative claims in this area.

\begin{thebibliography}{99}
\small

\bibitem{askin2026datamediated}
B.~Askin, M.~Ustaomeroglu, A.~Nayak, G.~Joshi, G.~Qu, and C.~Joe-Wong.
\newblock Emergent and Subliminal Misalignment Through the Lens of Data-Mediated Transfer.
\newblock arXiv:2605.12798, 2026.

\bibitem{betley2025em}
J.~Betley, D.~Tan, N.~Warncke, A.~Sztyber-Betley, X.~Bao, M.~Soto, N.~Labenz, and O.~Evans.
\newblock Emergent Misalignment: Narrow Finetuning Can Produce Broadly Misaligned LLMs.
\newblock arXiv:2502.17424, 2025.

\bibitem{betley2025weird}
J.~Betley, J.~Cocola, D.~Feng, J.~Chua, A.~Arditi, A.~Sztyber-Betley, and O.~Evans.
\newblock Weird Generalization and Inductive Backdoors: New Ways to Corrupt LLMs.
\newblock arXiv:2512.09742, 2025.

\bibitem{blank2026steering}
C.~Blank, A.~Bhatia, S.~Rajamanoharan, A.~Conmy, and N.~Nanda.
\newblock Subliminal Learning Is Steering Vector Distillation.
\newblock arXiv:2606.00995, 2026.

\bibitem{brockers2026noise}
V.~C. Brockers, R.~D. Ventzke, V.~Neuhaus, B.~Hidalgo-Ogalde, and V.~Priesemann.
\newblock Learning Through Noise: Why Subliminal Learning Works and When It Fails.
\newblock arXiv:2605.23645, 2026.

\bibitem{cho2026hidden}
M.~Cho, J.~Kim, S.~Song, J.~Yuk, M.~Kee, H.~Song, and K.~Lim.
\newblock Hidden Threat in Synthetic Data: Covert Targeted Bias Injection through Benign Text.
\newblock arXiv:2608.30619, 2026.

\bibitem{cloud2025subliminal}
A.~Cloud, M.~Le, J.~Chua, J.~Betley, A.~Sztyber-Betley, J.~Hilton, S.~Marks, and O.~Evans.
\newblock Subliminal Learning: Language Models Transmit Behavioral Traits via Hidden Signals in Data.
\newblock arXiv:2507.14805, 2025.

\bibitem{cloud2026nature}
A.~Cloud, M.~Le, J.~Chua, J.~Betley, A.~Sztyber-Betley, J.~Hilton, S.~Marks, and O.~Evans.
\newblock Language models transmit behavioural traits through hidden signals in data.
\newblock \emph{Nature}, 652:615, 2026. doi:10.1038/s41586-026-10319-8.

\bibitem{feng2026everyrank}
L.~Feng.
\newblock Subliminal Learning Happens at Every Rank, Given the Right Learning Rate and Enough Data.
\newblock LessWrong, 8 July 2026.

\bibitem{hadley2026nonsemantic}
E.~Hadley and E.~Gultepe.
\newblock Subliminal Learning is Non-Semantic Distillation.
\newblock arXiv:2608.05734, 2026.

\bibitem{hu2026salve}
N.~Hu, S.~Koyejo, and C.~Potts.
\newblock Verbalizing Subliminal Learning Effects Using Text Optimization.
\newblock arXiv:2609.16927, 2026.

\bibitem{liu2026drift}
Z.~Liu, Z.~Dong, Y.~Fan, X.~Li, and C.~Yang.
\newblock Subliminal Learning as Trait-Direction Drift: A Mechanism and Targeted Control under SFT Distillation.
\newblock arXiv:2609.01091, 2026.

\bibitem{morgulis2026steering}
G.~Morgulis and J.~Hewitt.
\newblock Subliminal Steering: Stronger Encoding of Hidden Signals.
\newblock arXiv:2604.25783, 2026.

\bibitem{narang2026consensus}
A.~Narang, A.~Tajdini, C.~Zhang, and J.~Morgenstern.
\newblock Inference-Time Consensus for Mitigating Hidden Behaviors from LLM Fine-Tuning.
\newblock arXiv:2607.23394, 2026.

\bibitem{nief2026lora}
T.~Nief, H.~Y. Fu, M.~Muchane, and A.~Holtzman.
\newblock Subliminal Learning is a LoRA Artifact.
\newblock arXiv:2606.00831, 2026.

\bibitem{schrodi2026understanding}
S.~Schrodi, E.~Kempf, F.~Barez, and T.~Brox.
\newblock Towards Understanding Subliminal Learning: When and How Hidden Biases Transfer.
\newblock ICLR 2026. arXiv:2509.23886.

\bibitem{soligo2025convergent}
A.~Soligo, E.~Turner, S.~Rajamanoharan, and N.~Nanda.
\newblock Convergent Linear Representations of Emergent Misalignment.
\newblock arXiv:2506.11618, 2025.

\bibitem{treutlein2024dots}
J.~Treutlein, D.~Choi, J.~Betley, S.~Marks, C.~Anil, R.~Grosse, and O.~Evans.
\newblock Connecting the Dots: LLMs Can Infer and Verbalize Latent Structure from Disparate Training Data.
\newblock arXiv:2406.14546, 2024.

\bibitem{turner2025organisms}
E.~Turner, A.~Soligo, M.~Taylor, S.~Rajamanoharan, and N.~Nanda.
\newblock Model Organisms for Emergent Misalignment.
\newblock arXiv:2506.11613, 2025.

\bibitem{vanderweijden2026repro}
D.~van~der Weijden, N.~Brack, and S.~Baez Santamaria.
\newblock Reproducing and Evaluating the Generalizability of Subliminal Learning in Open-Weight Models.
\newblock arXiv:2609.12586, 2026.

\bibitem{vir2025corruption}
R.~Vir and S.~Bhatnagar.
\newblock Subliminal Corruption: Mechanisms, Thresholds, and Interpretability.
\newblock arXiv:2510.19152, 2025.

\bibitem{wang2025oocr}
A.~Wang, J.~Engels, O.~Clive-Griffin, S.~Rajamanoharan, and N.~Nanda.
\newblock Simple Mechanistic Explanations for Out-of-Context Reasoning.
\newblock arXiv:2507.08218, 2025.

\bibitem{weckbecker2026thoughtvirus}
M.~Weckbecker, J.~M\"uller, B.~Hagag, and M.~Mulet.
\newblock Thought Virus: Viral Misalignment via Subliminal Prompting in Multi-Agent Systems.
\newblock arXiv:2603.00131, 2026.

\bibitem{zur2025owl}
A.~Zur, A.~Loftus, H.~Orgad, Z.~Ying, K.~Sahin, and D.~Bau.
\newblock It's Owl in the Numbers: Token Entanglement in Subliminal Learning.
\newblock Mechanistic Interpretability Workshop, NeurIPS 2025.

\end{thebibliography}

\end{document}
